% ==============================================================================
% CAPITOLO 18: MACROECONOMIA DEI SISTEMI ECONOMICI APERTI
% ==============================================================================
\chapter{Macroeconomia dei Sistemi Economici Aperti: Il Modello Mundell-Fleming}
\label{chap:macroeconomia_sistemi_aperti}

\begin{tcolorbox}[colback=navyblue!5!white, colframe=navyblue, title=\textbf{Quadro Concettuale del Capitolo}]
Questo capitolo conclude la trattazione della Macroeconomia sviluppando il celebre \textbf{Modello Mundell-Fleming} (IS-LM-BP per una piccola economia aperta) formulato da Robert Mundell e Marcus Fleming. Il modello integra il commercio internazionale di beni e la mobilità internazionale dei capitali finanziari nello schema IS-LM, determinando l'equilibrio congiunto del reddito, del tasso di interesse e del tasso di cambio. Verrà analizzata l'efficacia comparata della \textbf{politica fiscale} e della \textbf{politica monetaria} nei due regimi polari di \textbf{cambi fissi} e \textbf{cambi flessibili} in ipotesi di \textbf{perfetta mobilità dei capitali}. Si dimostrerà come in cambi fissi la moneta diventi endogena rendendo la politica fiscale massimamente efficace e quella monetaria impotente, mentre in cambi flessibili si verifichi l'opposto con lo \textbf{spiazzamento totale da tasso di cambio}. Si formalizzerà infine il \textbf{Trilemma Macroeconomico di Mundell-Obstfeld} (Trinità Impossibile).
\end{tcolorbox}

\section{Architettura del Modello Mundell-Fleming (IS-LM-BP)}

Il modello estende lo schema IS-LM a una \textbf{piccola economia aperta} inserita nei mercati internazionali dei beni e dei capitali.

\begin{modello}{Struttura Formale del Modello IS-LM-BP}{modello_mundell_fleming}
Il modello si compone di tre equazioni macroeconomiche fondamentali:
\begin{equation}
    \begin{cases}
        \IS: & Y = \alpha_G \Big[ A_0 - d \cdot r + NX(Y, \cambionom) \Big] \\[0.6em]
        \LM: & \dfrac{M}{P} = k Y - h r \\[0.8em]
        \BP: & NX(Y, \cambionom) + CF(r - r^*) = 0
    \end{cases} \label{eq:sistema_mf}
\end{equation}
dove:
\begin{itemize}
    \item $NX(Y, \cambionom) = X_0 - m Y - n \cambionom$ rappresenta le esportazioni nette (bilancia commerciale), con propensione marginale all'importazione $m = \pdv{M}{Y} > 0$ e sensibilità al cambio $n > 0$;
    \item $CF(r - r^*) = \sigma (r - r^*)$ rappresenta i movimenti netti di capitale finanziario privato, guidati dal differenziale tra il tasso di interesse interno $r$ e il tasso di interesse mondiale esogeno $r^*$, con parametro di mobilità $\sigma \ge 0$;
    \item $\cambionom$ è il tasso di cambio nominale quotato all'europea ($\text{USD/EUR}$).
\end{itemize}
\end{modello}

\subsection{La Curva BP e i Gradi di Mobilità Internazionale dei Capitali}

\begin{definizione}{Curva BP}{def_curva_bp}
La \textbf{Curva BP} è il luogo geometrico di tutte le combinazioni di reddito reale $Y$ e tasso di interesse $r$ che assicurano l'equilibrio complessivo della Bilancia dei Pagamenti ($\text{Saldo BP} = 0$, ovvero $\mathrm{CA} + \mathrm{FA}_{\text{priv}} = 0$).
\end{definizione}

Esplicitando il tasso di interesse $r$ dall'equazione della BP:
\begin{equation}
    r = r^* + \frac{m}{\sigma} Y - \frac{X_0 - n \cambionom}{\sigma}
\end{equation}
La pendenza della curva BP è $\dv{r}{Y}\big|_{\BP} = \frac{m}{\sigma} \ge 0$. Si distinguono tre scenari:
\begin{enumerate}
    \item \textbf{Assenza di Mobilità dei Capitali ($\sigma = 0$)}: i capitali non possono attraversare le frontiere; la curva BP è una retta \textbf{verticale} in corrispondenza del reddito che azzera il disavanzo commerciale: $Y = \frac{X_0 - n \cambionom}{m}$.
    \item \textbf{Mobilità Imperfetta dei Capitali ($0 < \sigma < \infty$)}: la curva BP è \textbf{inclinata positivamente}. Un aumento di reddito genera maggiori importazioni ($m \Delta Y > 0$); per ripristinare il pareggio della BP occorre un rialzo del tasso di interesse interno ($r > r^*$) che attiri un afflusso compensativo di capitali finanziari.
    \item \textbf{Perfetta Mobilità dei Capitali ($\sigma \to \infty$)}: le attività finanziarie nazionali ed estere sono perfette sostitute e i capitali si muovono istantaneamente al minimo differenziale di rendimento.
\end{enumerate}

\begin{teorema}{Curva BP con Perfetta Mobilità dei Capitali}{bp_perfetta_mobilita}
In condizioni di perfetta mobilità internazionale dei capitali ($\sigma \to \infty$), qualsiasi scostamento infinitesimale di $r$ da $r^*$ genera flussi di capitale di ampiezza illimitata. Pertanto, la condizione di equilibrio della bilancia dei pagamenti impone:
\begin{equation}
    r = r^* \label{eq:bp_orizzontale}
\end{equation}
La curva BP nel piano cartesiano $(Y, r)$ è una retta \textbf{perfettamente orizzontale} al livello del tasso di interesse internazionale $r^*$.
\end{teorema}

\section{Politiche Macroeconomiche in Regime di Cambi Fissi}

In regime di cambi fissi, il tasso di cambio è ancorato alla parità ufficiale $\bar{\cambionom}$ e la Banca Centrale è obbligata a intervenire acquistando o vendendo valuta estera ($\Delta RU \neq 0$).

\subsection{1. Politica Fiscale Espansiva: Massima Efficacia}

\begin{teorema}{Efficacia della Politica Fiscale in Cambi Fissi}{fiscale_cambi_fissi}
In un regime di cambi fissi con perfetta mobilità dei capitali, la politica fiscale ha la \textbf{massima efficacia espansiva}: non subisce alcuno spiazzamento finanziario e il PIL cresce dell'intero moltiplicatore dell'economia aperta.
\end{teorema}

\begin{dimostrazione}[ del Meccanismo di Trasmissione Fiscale in Cambi Fissi]
\begin{enumerate}
    \item \textbf{Shock Iniziale}: il governo aumenta la spesa pubblica di $\Delta G > 0$. La curva $\IS$ trasla verso destra da $\IS_0$ a $\IS_1$.
    \item \textbf{Pressione sui Tassi}: nel punto di intersezione provvisorio tra $\IS_1$ e $\LM_0$, il tasso di interesse interno sale al di sopra di quello mondiale: $r_1 > r^*$.
    \item \textbf{Afflusso di Capitali}: il differenziale positivo $(r_1 - r^* > 0)$ induce gli investitori internazionali a comprare massicciamente titoli nazionali, generando un enorme afflusso di valuta estera e una forte \textbf{pressione all'apprezzamento} del tasso di cambio.
    \item \textbf{Intervento Obbligato della Banca Centrale}: per impedire l'apprezzamento e difendere la parità fissa $\bar{\cambionom}$, la Banca Centrale è costretta ad acquistare la valuta estera offerta in eccesso, immettendo sul mercato moneta nazionale ($\Delta M > 0, \Delta RU > 0$).
    \item \textbf{Accomodamento Monetario Endogeno}: l'incremento automatico della base monetaria sposta la curva $\LM$ verso destra da $\LM_0$ a $\LM_1$.
    \item \textbf{Equilibrio Finale}: la traslazione della LM prosegue finché il tasso di interesse interno non ritorna esattamente al livello mondiale $r = r^*$.
    Il reddito di equilibrio finale è:
    \begin{equation}
        \Delta Y^* = \alpha_{\text{aperto}} \Delta G = \frac{1}{1 - c(1-t) + m} \Delta G > 0
    \end{equation}
\end{enumerate}
\end{dimostrazione}

\subsection{2. Politica Monetaria Espansiva: Totale Inefficacia}

\begin{teorema}{Inefficacia della Politica Monetaria in Cambi Fissi}{monetaria_cambi_fissi}
In un regime di cambi fissi con perfetta mobilità dei capitali, la Banca Centrale \textbf{perde completamente il controllo dell'offerta di moneta}: l'offerta di moneta diviene una variabile endogena e qualsiasi tentativo di espansione monetaria è totalmente inefficace sul PIL ($\Delta Y = 0$).
\end{teorema}

\begin{dimostrazione}[ dell'Impotenza Monetaria in Cambi Fissi]
\begin{enumerate}
    \item La Banca Centrale acquista titoli sul mercato aperto per immettere liquidità $\Delta M > 0$; la curva $\LM$ trasla temporaneamente a destra da $\LM_0$ a $\LM_1$.
    \item Il tasso di interesse interno scende al di sotto del livello internazionale: $r_1 < r^*$.
    \item I capitali finanziari fuggono istantaneamente all'estero alla ricerca di rendimenti più elevati: si genera un massiccio disavanzo della bilancia dei pagamenti e una fortissima \textbf{pressione al deprezzamento} della valuta nazionale.
    \item Per difendere la parità fissa $\bar{\cambionom}$, la Banca Centrale deve intervenire vendendo le proprie riserve ufficiali di valuta estera in cambio di moneta nazionale ($\Delta RU < 0$).
    \item Il riassorbimento della moneta nazionale distrugge base monetaria, costringendo la curva $\LM$ a \textbf{retrocedere verso sinistra da $\LM_1$ a $\LM_0$}.
    \item L'equilibrio finale coincide esattamente con quello di partenza: $Y^* = Y_0, r^* = r^*$, con l'unica conseguenza di una perdita secca di riserve valutarie della Banca Centrale.
\end{enumerate}
\end{dimostrazione}

\section{Politiche Macroeconomiche in Regime di Cambi Flessibili}

In regime di cambi flessibili, il tasso di cambio oscilla liberamente sul Forex per garantire l'equilibrio della bilancia dei pagamenti ($\Delta RU = 0$).

\subsection{1. Politica Fiscale Espansiva: Spiazzamento Totale da Cambio}

\begin{teorema}{Inefficacia della Politica Fiscale in Cambi Flessibili}{fiscale_cambi_flessibili}
In un regime di cambi flessibili con perfetta mobilità dei capitali, la politica fiscale è \textbf{completamente inefficace nel modificare il livello del PIL} ($\Delta Y = 0$), a causa del fenomeno dello \textbf{Spiazzamento Totale da Tasso di Cambio}.
\end{teorema}

\begin{dimostrazione}[ dello Spiazzamento da Cambio]
\begin{enumerate}
    \item Il governo incrementa la spesa pubblica di $\Delta G > 0$; la curva $\IS$ si sposta a destra da $\IS_0$ a $\IS_1$.
    \item Il tasso di interesse interno tende a salire ($r_1 > r^*$), attirando ingenti flussi di capitale estero.
    \item In assenza di interventi della Banca Centrale, l'eccesso di domanda di valuta nazionale provoca un \textbf{apprezzamento immediato del tasso di cambio nominale} ($\cambionom \uparrow$).
    \item L'apprezzamento rende i beni nazionali più costosi per gli acquirenti esteri e i beni esteri più convenienti per i residenti: la competitività crolla e le \textbf{esportazioni nette si riducono} ($\Delta NX < 0$).
    \item La contrazione di $NX$ sposta la curva $\IS$ all'indietro verso sinistra finché essa non ritorna esattamente nella posizione iniziale $\IS_0$.
    \item L'equilibrio finale si ricompone al reddito iniziale $Y_0$ al tasso $r^*$: la spesa pubblica ha spiazzato le esportazioni nette in misura esattamente equivalente ($\Delta NX = -\Delta G$), lasciando il PIL inalterato ($\Delta Y = 0$).
\end{enumerate}
\end{dimostrazione}

\subsection{2. Politica Monetaria Espansiva: Massima Efficacia}

\begin{teorema}{Efficacia della Politica Monetaria in Cambi Flessibili}{monetaria_cambi_flessibili}
In un regime di cambi flessibili con perfetta mobilità dei capitali, la politica monetaria acquisisce la \textbf{massima efficacia espansiva}: la trasmissione non passa attraverso la riduzione del tasso di interesse (che resta vincolato a $r = r^*$), bensì attraverso il \textbf{canale del deprezzamento del tasso di cambio e l'espansione delle esportazioni nette}.
\end{teorema}

\begin{dimostrazione}[ della Trasmissione Monetaria via Tasso di Cambio]
\begin{enumerate}
    \item La Banca Centrale aumenta l'offerta di moneta $\Delta M > 0$; la curva $\LM$ trasla verso destra da $\LM_0$ a $\LM_1$.
    \item Il tasso interno scende temporaneamente al di sotto di quello mondiale ($r_1 < r^*$), innescando un immediato deflusso di capitali finanziari verso l'estero.
    \item La vendita di moneta nazionale sul mercato Forex provoca un \textbf{forte deprezzamento del tasso di cambio} ($\cambionom \downarrow$).
    \item Il deprezzamento aumenta la competitività di prezzo delle merci nazionali: le esportazioni aumentano e le importazioni calano ($\Delta NX > 0$).
    \item L'incremento della domanda estera sposta la curva $\IS$ verso destra da $\IS_0$ a $\IS_1$, finché non interseca la nuova $\LM_1$ esattamente in corrispondenza del tasso mondiale $r = r^*$.
    \item Il PIL di equilibrio aumenta pienamente:
    \begin{equation}
        \Delta Y^* = \frac{1}{k}\Delta\left(\frac{M}{P}\right) > 0
    \end{equation}
\end{enumerate}
\end{dimostrazione}

\begin{figure}[htbp]
\centering
\begin{tikzpicture}[scale=0.95, >=Stealth]
    % GRAFICO A SINISTRA: CAMBI FISSI
    \begin{scope}[xshift=0cm]
        \draw[->, line width=1.1pt] (0,0) -- (6.8,0) node[right, font=\small\bfseries] {$Y$};
        \draw[->, line width=1.1pt] (0,0) -- (0,5.5) node[above, font=\small\bfseries] {$r$};
        \node[font=\footnotesize\bfseries, color=navyblue] at (3.4, 5.2) {(a) Cambi Fissi: Fiscale Espansiva};

        % BP Orizzontale
        \draw[line width=1.8pt, color=forestgreen] (0, 2.5) -- (6.2, 2.5) node[right, font=\footnotesize\bfseries] {$\mathrm{BP} \; (r = r^*)$};
        \node[left, font=\footnotesize\bfseries, color=forestgreen] at (0, 2.5) {$r^*$};

        % Curve IS
        \draw[line width=1.5pt, color=navyblue] (0.8, 4.5) -- (4.8, 0.5) node[right, font=\scriptsize\bfseries] {$\IS_0$};
        \draw[line width=1.6pt, color=navyblue!60!blue] (2.2, 5.2) -- (6.2, 1.2) node[right, font=\scriptsize\bfseries] {$\IS_1 (\Delta G > 0)$};

        % Curve LM
        \draw[line width=1.5pt, color=purpleaccent] (0.8, 0.5) -- (4.8, 4.5) node[above left, font=\scriptsize\bfseries] {$\LM_0$};
        \draw[line width=1.6pt, dashed, color=tealblue] (2.2, 0.5) -- (6.2, 4.5) node[right, font=\scriptsize\bfseries] {$\LM_1 (\text{Endogena})$};

        % Equilibri
        \coordinate (E0) at (2.8, 2.5);
        \coordinate (E1) at (4.2, 2.5);

        \fill[navyblue] (E0) circle (2.5pt) node[above left=2pt, font=\scriptsize\bfseries] {$E_0$};
        \fill[forestgreen] (E1) circle (2.5pt) node[above right=2pt, font=\scriptsize\bfseries] {$E_1$};

        \draw[dotted, line width=0.8pt, color=navyblue] (2.8, 0) node[below, font=\scriptsize\bfseries] {$Y_0^*$} -- (E0);
        \draw[dotted, line width=0.8pt, color=forestgreen] (4.2, 0) node[below, font=\scriptsize\bfseries] {$Y_1^*$} -- (E1);
        \draw[->, line width=1.2pt, color=crimsonred] (2.8, 2.7) -- node[above, font=\tiny\bfseries] {$\Delta Y > 0$} (4.2, 2.7);
    \end{scope}

    % GRAFICO A DESTRA: CAMBI FLESSIBILI
    \begin{scope}[xshift=7.8cm]
        \draw[->, line width=1.1pt] (0,0) -- (6.8,0) node[right, font=\small\bfseries] {$Y$};
        \draw[->, line width=1.1pt] (0,0) -- (0,5.5) node[above, font=\small\bfseries] {$r$};
        \node[font=\footnotesize\bfseries, color=purpleaccent] at (3.4, 5.2) {(b) Cambi Flessibili: Monetaria Espansiva};

        % BP Orizzontale
        \draw[line width=1.8pt, color=forestgreen] (0, 2.5) -- (6.2, 2.5) node[right, font=\footnotesize\bfseries] {$\mathrm{BP} \; (r = r^*)$};
        \node[left, font=\footnotesize\bfseries, color=forestgreen] at (0, 2.5) {$r^*$};

        % Curve LM
        \draw[line width=1.5pt, color=purpleaccent] (0.8, 0.5) -- (4.8, 4.5) node[above left, font=\scriptsize\bfseries] {$\LM_0$};
        \draw[line width=1.6pt, color=purpleaccent!60!black] (2.2, 0.5) -- (6.2, 4.5) node[right, font=\scriptsize\bfseries] {$\LM_1 (\Delta M > 0)$};

        % Curve IS
        \draw[line width=1.5pt, color=navyblue] (0.8, 4.5) -- (4.8, 0.5) node[right, font=\scriptsize\bfseries] {$\IS_0$};
        \draw[line width=1.6pt, dashed, color=tealblue] (2.2, 5.2) -- (6.2, 1.2) node[right, font=\scriptsize\bfseries] {$\IS_1 (e \downarrow, NX \uparrow)$};

        % Equilibri
        \coordinate (E0f) at (2.8, 2.5);
        \coordinate (E1f) at (4.2, 2.5);

        \fill[purpleaccent] (E0f) circle (2.5pt) node[above left=2pt, font=\scriptsize\bfseries] {$E_0$};
        \fill[forestgreen] (E1f) circle (2.5pt) node[above right=2pt, font=\scriptsize\bfseries] {$E_1$};

        \draw[dotted, line width=0.8pt, color=purpleaccent] (2.8, 0) node[below, font=\scriptsize\bfseries] {$Y_0^*$} -- (E0f);
        \draw[dotted, line width=0.8pt, color=forestgreen] (4.2, 0) node[below, font=\scriptsize\bfseries] {$Y_1^*$} -- (E1f);
        \draw[->, line width=1.2pt, color=crimsonred] (2.8, 2.7) -- node[above, font=\tiny\bfseries] {$\Delta Y > 0$} (4.2, 2.7);
    \end{scope}

\end{tikzpicture}
\caption{Efficacia comparata nel Modello Mundell-Fleming con perfetta mobilità dei capitali: (a) In Cambi Fissi la politica fiscale è massimamente efficace grazie all'accomodamento monetario endogeno; (b) In Cambi Flessibili la politica monetaria è massimamente efficace grazie al deprezzamento del cambio e al traino delle esportazioni nette.}
\label{fig:mundell_fleming_cambi_fissi_flessibili}
\end{figure}

\subsection{Quadro Sinottico Comparativo di Mundell-Fleming}

\begin{table}[htbp]
\centering
\small
\renewcommand{\arraystretch}{1.4}
\begin{tabularx}{\textwidth}{l X X}
\toprule
\textbf{Regime di Cambio} & \textbf{Politica Fiscale ($\Delta G > 0$)} & \textbf{Politica Monetaria ($\Delta M > 0$)} \\
\midrule
\textbf{Cambi Fissi} & 
\textbf{Massimamente Efficace}\newline
$\Delta Y = \alpha_{\text{aperto}} \Delta G$\newline
($r = r^*$, $\LM$ si sposta a destra, zero spiazzamento) & 
\textbf{Completamente Inefficace}\newline
$\Delta Y = 0$\newline
(la moneta è endogena; perdita di riserve valutarie) \\
\midrule
\textbf{Cambi Flessibili} & 
\textbf{Completamente Inefficace}\newline
$\Delta Y = 0$\newline
(spiazzamento totale da cambio: $\cambionom \uparrow \implies \Delta NX = -\Delta G$) & 
\textbf{Massimamente Efficace}\newline
$\Delta Y = \frac{1}{k}\Delta \frac{M}{P}$\newline
(deprezzamento del cambio: $\cambionom \downarrow \implies NX \uparrow$) \\
\bottomrule
\end{tabularx}
\caption{Matrice riassuntiva dell'efficacia delle politiche economiche nel Modello Mundell-Fleming.}
\label{tab:sinottico_mundell_fleming}
\end{table}

\section{Il Trilemma Macroeconomico (Trinità Impossibile di Mundell)}

\begin{teorema}{Il Trilemma Macroeconomico (Mundell-Obstfeld)}{trilemma_mundell}
Un'economia nazionale integrata nel sistema mondiale può scegliere \textbf{contemporaneamente al massimo due} dei seguenti tre regimi istituzionali:
\begin{enumerate}
    \item \textbf{Stabilità del Tasso di Cambio (Regime di Cambi Fissi)};
    \item \textbf{Perfetta Mobilità Internazionale dei Capitali (Integrazione Finanziaria)};
    \item \textbf{Autonomia della Politica Monetaria Nazionale}.
\end{enumerate}
È analiticamente e operativamente impossibile conseguire simultaneamente tutti e tre gli obiettivi.
\end{teorema}

\begin{figure}[htbp]
\centering
\begin{tikzpicture}[scale=1.1, >=Stealth]
    % Vertici del Triangolo
    \coordinate (V1) at (0, 0);          % Cambi Fissi
    \coordinate (V2) at (7.0, 0);        % Autonomia Monetaria
    \coordinate (V3) at (3.5, 5.5);      % Perfetta Mobilità Capitali

    % Disegno lati del triangolo
    \draw[line width=2pt, color=navyblue] (V1) -- (V2) -- (V3) -- cycle;

    % Nodi ai vertici
    \node[rectangle, draw=navyblue, fill=navyblue!12!white, line width=1.2pt, rounded corners=4pt, align=center, font=\small\bfseries] at (V1) {
        \textbf{1. Cambi Fissi}\\
        (Stabilità Valutaria)
    };

    \node[rectangle, draw=purpleaccent, fill=purpleaccent!12!white, line width=1.2pt, rounded corners=4pt, align=center, font=\small\bfseries] at (V2) {
        \textbf{2. Politica Monetaria}\\
        \textbf{Autonoma}
    };

    \node[rectangle, draw=forestgreen, fill=forestgreen!12!white, line width=1.2pt, rounded corners=4pt, align=center, font=\small\bfseries] at (V3) {
        \textbf{3. Libera Circolazione}\\
        \textbf{dei Capitali}
    };

    % Etichette sui tre lati con esempi storici
    \node[align=center, font=\scriptsize\bfseries, color=navyblue, rotate=58] at (1.5, 3.0) {
        \textbf{Opzione A: Eurozona / Gold Standard}\\
        (Cambi Fissi + Libera Mobilità)\\
        $\implies$ Rinuncia alla sovranità monetaria
    };

    \node[align=center, font=\scriptsize\bfseries, color=purpleaccent, rotate=-58] at (5.5, 3.0) {
        \textbf{Opzione B: USA / UK / Giappone}\\
        (Cambi Flessibili + Libera Mobilità)\\
        $\implies$ Rinuncia alla stabilità del cambio
    };

    \node[align=center, font=\scriptsize\bfseries, color=crimsonred] at (3.5, -0.6) {
        \textbf{Opzione C: Accordi di Bretton Woods (1944-1971) / Cina anni '90}\\
        (Cambi Fissi + Politica Monetaria Autonoma)\\
        $\implies$ Rinuncia alla libera mobilità dei capitali (Controlli sui capitali)
    };

\end{tikzpicture}
\caption{La Trinità Impossibile (Trilemma di Mundell). Ciascun Paese deve collocarsi necessariamente su uno dei tre lati del triangolo, rinunciando all'obiettivo posto al vertice opposto.}
\label{fig:trilemma_mundell}
\end{figure}

\section{Metodi Risolutivi ed Eserciziario d'Esame}

\begin{metodo}[Algoritmo Risolutivo per il Modello Mundell-Fleming]
Per risolvere un problema di economia aperta con perfetta mobilità dei capitali:
\begin{enumerate}
    \item \textbf{Identificazione del Regime di Cambio}:
    \begin{itemize}
        \item Se \textbf{Cambi Fissi}: fissare $r = r^*$, $\cambionom = \bar{\cambionom}$. Il reddito $Y^*$ è determinato interamente dalla curva IS al tasso $r^*$. L'offerta di moneta $M$ è endogena e si ricava sostituendo $(Y^*, r^*)$ nella LM.
        \item Se \textbf{Cambi Flessibili}: fissare $r = r^*$ e offerta di moneta $M = \bar{M}$. Il reddito $Y^*$ è determinato interamente dalla curva LM al tasso $r^*$. Il tasso di cambio nominale di equilibrio $\cambionom^*$ si ricava sostituendo $(Y^*, r^*)$ nella curva IS.
    \end{itemize}
    \item \textbf{Valutazione degli Shock}:
    \begin{itemize}
        \item In cambi fissi con $\Delta G$: calcolare $\Delta Y = \alpha_{\text{aperto}}\Delta G$ e l'incremento di moneta $\Delta M = P(k \Delta Y)$.
        \item In cambi flessibili con $\Delta G$: $\Delta Y = 0$; calcolare la svalutazione/apprezzamento necessario ad assorbire lo shock: $\Delta NX = -\Delta G$.
    \end{itemize}
\end{enumerate}
\end{metodo}

\begin{esercizio}[Risoluzione Completa Mundell-Fleming in Cambi Fissi e Flessibili]
Si consideri una piccola economia aperta descritta dalle seguenti relazioni:
\begin{align*}
    C &= 200 + 0{,}8(Y - T), \quad T = 0{,}25 Y, \quad I = 400 - 2000 r, \quad G = 300 \\
    NX &= 150 - 0{,}1 Y - 100 \cambionom \\
    \left(\frac{M}{P}\right)^{\mathrm{d}} &= 0{,}5 Y - 1000 r, \quad P = 1{,}0, \quad r^* = 0{,}05 \; (5\%)
\end{align*}
\textbf{Richieste}:
\begin{enumerate}
    \item Determinare l'equazione della curva IS ed esprimere $Y$ in funzione di $r$ ed $\cambionom$.
    \item \textbf{Regime di Cambi Fissi} con parità fissata a $\bar{\cambionom} = 1{,}0$: calcolare il reddito di equilibrio $Y^*$, lo stock di moneta di equilibrio $M^*$, il consumo $C^*$, l'investimento $I^*$ e il saldo commerciale $NX^*$.
    \item Valutare l'effetto di un incremento di spesa pubblica $\Delta G = 50$ in regime di cambi fissi.
    \item \textbf{Regime di Cambi Flessibili} con stock di moneta fissato esogenamente a $\bar{M} = 850$: calcolare il reddito di equilibrio $Y^*$ e il tasso di cambio di equilibrio $\cambionom^*$. Valutare l'effetto di $\Delta G = 50$.
\end{enumerate}

\textbf{Svolgimento}:
\begin{enumerate}
    \item \textbf{Equazione della Curva IS}:
    \begin{align*}
        Y &= C + I + G + NX = 200 + 0{,}8(0{,}75 Y) + 400 - 2000 r + 300 + 150 - 0{,}1 Y - 100 \cambionom \\
        Y &= 1050 + 0{,}6 Y - 0{,}1 Y - 2000 r - 100 \cambionom = 1050 + 0{,}5 Y - 2000 r - 100 \cambionom \\
        0{,}5 Y &= 1050 - 2000 r - 100 \cambionom \implies Y = 2100 - 4000 r - 200 \cambionom \quad (\alpha_{\text{aperto}} = 2{,}0)
    \end{align*}

    \item \textbf{Cambi Fissi ($\bar{\cambionom} = 1{,}0, r = r^* = 0{,}05$)}:
    \begin{align*}
        Y^* &= 2100 - 4000(0{,}05) - 200(1{,}0) = 2100 - 200 - 200 = 1700
    \end{align*}
    Dalla curva LM ricaviamo l'offerta endogena di moneta $M^*$:
    \begin{equation*}
        \frac{M^*}{1{,}0} = 0{,}5(1700) - 1000(0{,}05) = 850 - 50 = 800 \implies M^* = 800
    \end{equation*}
    Componenti dell'equilibrio:
    \begin{align*}
        C^* &= 200 + 0{,}8(0{,}75 \times 1700) = 200 + 0{,}6(1700) = 200 + 1020 = 1220 \\
        I^* &= 400 - 2000(0{,}05) = 300 \\
        NX^* &= 150 - 0{,}1(1700) - 100(1{,}0) = 150 - 170 - 100 = -120 \quad (\text{Disavanzo})
    \end{align*}
    Verifica: $C^* + I^* + G + NX^* = 1220 + 300 + 300 - 120 = 1700 = Y^*$.

    \item \textbf{Effetto di $\Delta G = 50$ in Cambi Fissi}:
    Nuova spesa autonoma: $1050 + 50 = 1100 \implies \text{Nuova IS}: Y = 2200 - 4000 r - 200 \cambionom$.
    Nuovo reddito:
    \begin{equation*}
        Y_1^* = 2200 - 200 - 200 = 1800 \implies \Delta Y = 1800 - 1700 = +100 = \alpha_{\text{aperto}} \Delta G = 2 \times 50
    \end{equation*}
    Nuovo stock di moneta endogeno:
    \begin{equation*}
        M_1^* = 0{,}5(1800) - 50 = 900 - 50 = 850 \implies \Delta M = +50
    \end{equation*}
    La politica fiscale è pienamente efficace con moltiplicatore $2{,}0$.

    \item \textbf{Cambi Flessibili con $\bar{M} = 850$}:
    In cambi flessibili il reddito è bloccato dalla LM al tasso $r = r^* = 0{,}05$:
    \begin{equation*}
        850 = 0{,}5 Y - 1000(0{,}05) \implies 0{,}5 Y = 850 + 50 = 900 \implies Y^* = 1800
    \end{equation*}
    Sostituendo $Y = 1800$ e $r = 0{,}05$ nell'equazione della IS iniziale:
    \begin{equation*}
        1800 = 2100 - 4000(0{,}05) - 200 \cambionom = 1900 - 200 \cambionom \implies 200 \cambionom = 100 \implies \cambionom^* = 0{,}50
    \end{equation*}
    Se ora il governo aumenta la spesa di $\Delta G = 50$, dalla LM il reddito rimane inchiodato a $Y^* = 1800$ ($\Delta Y = 0$).
    Dalla nuova equazione IS ($Y = 2200 - 4000 r - 200 \cambionom$):
    \begin{equation*}
        1800 = 2200 - 200 - 200 \cambionom_1 \implies 200 \cambionom_1 = 200 \implies \cambionom_1^* = 1{,}00
    \end{equation*}
    Il tasso di cambio si apprezza da $0{,}50$ a $1{,}00$ (+100 per cento), provocando una contrazione delle esportazioni nette esattamente di:
    \begin{equation*}
        \Delta NX = -100(1{,}00 - 0{,}50) = -50 = -\Delta G \quad (\textbf{Spiazzamento Totale da Cambio})
    \end{equation*}
\end{enumerate}
\end{esercizio}

\begin{esame}[Tranelli Cruciali sul Modello Mundell-Fleming]
\begin{enumerate}
    \item \textbf{Spiazzamento Fiscale}: in cambi fissi non c'è spiazzamento (massima efficacia); in cambi flessibili lo spiazzamento è del 100 per cento, ma avviene sul \textbf{tasso di cambio reale e sulle esportazioni nette} ($NX \downarrow$), non sul tasso di interesse o sugli investimenti ($r = r^*$).
    \item \textbf{Sovranità Monetaria}: all'esame si chiede spesso perché un Paese dell'Eurozona non possa attuare una politica monetaria nazionale. La risposta risiede nel Trilemma: avendo cambi fissi irrevocabili (euro) e perfetta mobilità dei capitali, la Banca Centrale nazionale non può manovrare autonomamente l'offerta di moneta o i tassi di interesse.
\end{enumerate}
\end{esame}
